% Part: lambda-calculus
% Chapter: introduction
% Section: primitive-recursion

\documentclass[../../../include/open-logic-section]{subfiles}

\begin{document}

\olfileid{lam}{int}{pr}
\olsection{\usetoken{S}{lambda definable} ప్రమేయాలు ఆదిమ పునరావృత్తి కింద సంవృతమైనవి}

ఆదిమ పునరావృత్తి విషయానికి వచ్చేసరికి, చివరకు మనం కొంత పని చేయాలి.
దశలవారీగా ముందుకు సాగాలి. మునుపటిలాగే, $G$, $H$లు వరుసగా $g$,~$h$
ప్రమేయాలను \tetoken{లాంబ్డాతో నిర్వచించే}{!!{lambda define}} పదాలు ఇప్పటికే ఉన్నాయని అనుకుందాం. కింది
విధంగా నిర్వచించిన ప్రమేయం~$f$కు $F$ అనే పదం కావాలి; అది దానిని
\tetoken{లాంబ్డాతో నిర్వచిస్తుంది}{!!{lambda define}s}:
\begin{align*}
f(0, \vec z) & = g(\vec z) \\
f(x+1, \vec z) & = h(x, f(x,\vec z), \vec z).
\end{align*}
\sourcecorrection{OLTELAMCOMP-003}{ఆంగ్ల మూలం దశ ప్రమేయం~$h$ను నిర్వచించే పదానికి ఇప్పటికే $H$ను వాడి, తరువాత $f$ను నిర్వచించాల్సిన పదానికీ $H$ అనే పేరే అడిగింది. వెంటనే వచ్చే నిర్మాణమంతా వాడే ఉద్దేశిత ఫలిత పదం~$F$ను పునరుద్ధరించాం.}
\sourcecorrection{OLTELAMCOMP-004}{ఆదిమ పునరావృత్తి దశలో మొదటి ఆర్గ్యుమెంటు పునరావృత్తి సూచిక~$x$ కావాలి; కానీ ఆంగ్ల మూలం అదనపు పరామితుల సదిశం~$\vec z$లోని $z$ను రాసింది. తరువాతి ప్రతి సమీకరణకు అనుగుణంగా $h(x,f(x,\vec z),\vec z)$ను పునరుద్ధరించాం.}
కాబట్టి సాధారణంగా, లాంబ్డా పదాలు $G'$,~$H'$ ఇచ్చినప్పుడు, ప్రతి సహజ
సంఖ్య~$n$కు కింది సమీకరణాలను నెరవేర్చే పదం~$F$ను కనుగొంటే సరిపోతుంది:
\begin{align*}
F(\num{0}, \vec z) & \equiv G(\vec z) \\
F(\overline{n+1}, \vec z) & \equiv H(\num{n}, F(\num{n}, \vec z), \vec z)
\end{align*}
$G'$,~$H'$లు $g$,~$h$లను \tetoken{లాంబ్డాతో నిర్వచించే}{!!{lambda define}} పదాలు కావడం వల్ల,
$\vec z$ స్థానంలో సంఖ్యాంకాలు $\num{\vec m}$ను పెట్టినప్పుడల్లా
$F(\num{n+1}, \num{\vec m})$ సరైన సమాధానానికి నియతీకృతమవుతుంది.

కానీ ఇందుకు ప్రతి సహజ సంఖ్య~$n$కు కింది సమీకరణాలను నెరవేర్చే పదం~$F$ను
కనుగొంటే సరిపోతుంది:
\begin{align*}
F(\num 0) & \equiv G \\
F(\overline {n+1}) & \equiv H(\num{n},F(\num{n}))
\intertext{ఇక్కడ}
G & = \lambd[\vec z][G'(\vec z)] \text{ మరియు}\\
H(u,v) & = \lambd[\vec z][H'(u,v(\vec z),\vec z)].
\end{align*}
\sourcecorrection{OLTELAMCOMP-005}{అదనపు పరామితులను అమూర్తీకరణలో ఇమిడ్చిన తరువాత $v$ అనేది $F(n)$గా, కేవలం $\vec z$కు ప్రమేయంగా ఉంటుంది. ఆంగ్ల మూలం దానికి పునరావృత్తి సూచిక~$u$ను మళ్లీ అదనపు ఆర్గ్యుమెంటుగా ఇచ్చింది; ఉద్దేశిత $v(\vec z)$ను పునరుద్ధరించాం.}
మరో మాటలో, లాంబ్డా చాతుర్యంతో అదనపు పరామితులు~$\vec z$ గురించి
వేరుగా ఆందోళన చెందనవసరం లేదు---అవి లాంబ్డా సంకేతనంలో ఇమిడిపోతాయి.

$F$ పదాన్ని నిర్వచించే ముందు, క్రమయుగ్మాలను నిర్వహించే యంత్రాంగం
అవసరం. తరువాతి ఉపసిద్ధాంతం దాన్ని అందిస్తుంది.

\begin{lem}
లాంబ్డా పదాల ప్రతి జత $M$,~$N$కు $D(M,N)(\num{0}) \red M$ మరియు
$D(M,N)(\num{1}) \red N$ అయ్యేలా ఒక లాంబ్డా పదం~$D$ ఉంది.
\end{lem}

\begin{proof}
మొదట, లాంబ్డా పదం~$K$ను కింది విధంగా నిర్వచించండి:
\[
K(y) = \lambd[x][y].
\]
మరో మాటలో, $K$ అనేది $\lambd[y][\lambd[x][y]]$ అనే పదం. వేరొక
విధంగా చూస్తే, ప్రతి~$M$కు $K(M)$ అనేది ఏ ఇన్‌పుట్‌పైనైనా~$M$ను
తిరిగి ఇచ్చే స్థిర ప్రమేయం.

ఇప్పుడు $D(x,y,z)$ను $D(x,y,z) = z (K(y))x$ ద్వారా నిర్వచించండి.
అప్పుడు మనకు కావలసినట్లుగానే
\begin{align*}
D(M,N,\num 0) & \red \num 0 (K(N)) M \red M \text{ మరియు}\\
D(M,N,\num 1) & \red \num 1 (K(N)) M \red K(N) M \red N,
\end{align*}
వస్తాయి.
\end{proof}

$D(M,N)$ అనేది $\tuple{M, N}$ జతకు ప్రాతినిధ్యం వహిస్తుందనేది భావన.
$P$ అలాంటి జతకు ప్రాతినిధ్యం వహిస్తుందని అనుకుంటే, $P(\num 0)$,
$P(\num 1)$లు ఎడమ, కుడి ప్రక్షేపాలు $(P)_0$, $(P)_1$లకు
ప్రాతినిధ్యం వహిస్తాయి. ఇకపై రెండవ సంకేతనాలనే వాడతాం.

\begin{lem}
\tetoken{లాంబ్డాతో నిర్వచించదగిన}{!!{lambda definable}} ప్రమేయాలు ఆదిమ పునరావృత్తి కింద సంవృతమైనవి.
\end{lem}

\begin{proof}
ఏ పదాలు $G$,~$H$ ఇచ్చినా, ప్రతి సహజ సంఖ్య~$n$కు
\begin{align*}
F(\num{0}) & \equiv G \\
F(\num{n+1}) & \equiv H(\num{n}, F(\num{n}))
\end{align*}
అయ్యే పదం~$F$ను కనుగొనగలమని చూపాలి. సంఖ్యాంకాలను పునరావర్తకాలుగా
వాడి, స్థూలంగా చెప్పాలంటే, ఈ
\emph{జతల} అనుక్రమాన్ని గణించాలనేది భావన:
\[
\tuple{\num 0, F(\num 0)}, \tuple{\num 1, F(\num 1)}, \dots,
\]
మొదటి జత కేవలం $\tuple{\num 0, G}$ అని గమనించండి. ఒక జత
$\tuple{\num{n}, F(\num{n})}$ ఇచ్చినప్పుడు, తరువాతి జత
$\tuple{\num{n+1}, F(\num{n+1})}$ అనేది
$\tuple{\num{n+1}, H(\num{n}, F(\num{n}))}$కు తుల్యంగా ఉండాలి.
ఈ ఒక-దశ మార్పును చేసే లాంబ్డా పదం~$T$ను రూపొందిస్తాం.

వివరాలు ఇవి. $T(u)$ను కింది విధంగా నిర్వచించండి:
\[
T(u) = \tuple{\fn{Succ}((u)_0), H((u)_0,(u)_1)}.
\]
\sourcecorrection{OLTELAMCOMP-006}{మునుపటి విభాగం ఉత్తరగామి లాంబ్డా పదాన్ని $\fn{Succ}$గా నిర్వచించినప్పటికీ, ఆంగ్ల మూలం ఇక్కడ నిర్వచించని $S$ను వాడింది. తరువాతి పంక్తిలో మొదటి భాగం $\num{n+1}$కు తగ్గాలనే ధృవీకరణకు అనుగుణంగా $\fn{Succ}$ను వాడాం.}
ఇప్పుడు ఏ సంఖ్య~$n$కైనా కింది సమీకరణాన్ని సులభంగా ధృవీకరించవచ్చు:
\[
T(\tuple{\num{n}, M}) \red \tuple{\num{n+1}, H(\num{n}, M)}.
\]
పైన సూచించినట్లుగా, $G$,~$H$ ఇచ్చినప్పుడు $F(u)$ను కింది విధంగా
నిర్వచించండి:
\[
F(u) = (u(T,\tuple{\num 0, G}))_1.
\]
మరో మాటలో, ఇన్‌పుట్~$\num{n}$పై $F$ అనేది $T$ను $\tuple{\num 0, G}$పై
$n$ సార్లు పునరావర్తించి, తరువాత రెండవ అంశాన్ని తిరిగి ఇస్తుంది. మొదటగా,
\begin{enumerate}
\item $\num{0} (T, \tuple{\num 0, G}) \equiv \tuple{\num{0}, G}$
\item $F(\num{0}) \equiv G$
\end{enumerate}
$n$పై ఆగమనం ద్వారా, ప్రతి సహజ సంఖ్యకు కిందివి ఉంటాయని చూపవచ్చు:
\begin{enumerate}
\item $\num{n+1}(T, \tuple{\num{0}, G}) \equiv \tuple{\num{n+1}, F(\num{n+1})}$
\item $F(\num{n+1}) \equiv H(\num{n}, F(\num{n}))$
\end{enumerate}
రెండవ ఉపవాక్యానికి,
\begin{align*}
F(\num{n+1}) & \red (\num{n+1}(T, \tuple{\num 0, G}))_1 \\
& \equiv (T(\num{n} (T, \tuple{\num 0, G})))_1 \\
& \equiv (T(\tuple{\num{n}, F(\num{n})}))_1 \\
& \equiv (\tuple{\num{n+1}, H(\num{n}, F(\num{n}))})_1 \\
& \equiv H(\num{n}, F(\num{n})).
\end{align*}
ఇక్కడ చివరి నుంచి రెండవ పంక్తిలో ఆగమన పరికల్పనను వాడాం. మొదటి
ఉపవాక్యానికి,
\begin{align*}
\num{n+1} (T, \tuple{\num 0, G}) &
\equiv T(\num{n} (T, \tuple{\num 0, G})) \\
& \equiv T( \tuple{\num{n}, F(\num{n})}) \\
& \equiv \tuple{\num{n+1}, H(\num{n}, F(\num{n}))} \\
& \equiv \tuple{\num{n+1}, F(\num{n+1})}.
\end{align*}
ఇక్కడ చివరి పంక్తిలో రెండవ ఉపవాక్యాన్ని వాడాం. ఈ విధంగా
$F(\num 0) \equiv G$ అనీ, ప్రతి~$n$కు $F(\num {n+1}) \equiv H(\num
n, F(\num{n}))$ అనీ చూపాం; మనకు కావలసింది సరిగ్గా ఇదే.
\end{proof}

\end{document}
